\documentclass[11pt]{article}
\usepackage[T1]{fontenc}
\usepackage{lmodern}
\usepackage[margin=1in]{geometry}
\usepackage{amsmath,amssymb,amsthm,booktabs,array}
\usepackage[colorlinks=true,urlcolor=blue,linkcolor=blue,citecolor=blue]{hyperref}
\hypersetup{pdftitle={A direct split-prime bound for the tenth discriminant: expanded v3},
 pdfauthor={Research note prepared with Codex},
 pdfsubject={Version 0.3.1: the bound q < 1,000,000 and its finite verification}}
\usepackage{microtype}
\newtheorem{theorem}{Theorem}
\newtheorem{lemma}[theorem]{Lemma}
\newtheorem{proposition}[theorem]{Proposition}
\newcommand{\Z}{\mathbb Z}
\newcommand{\Q}{\mathbb Q}
\newcommand{\OK}{\mathcal O_K}
\title{A direct split-prime bound for the tenth discriminant\\
\large Version 3, expanded exposition: \(q<10^6\)}
\author{Research note prepared with Codex}
\date{Version 0.3.1 --- 20 October 2026}
\begin{document}
\maketitle
\begin{abstract}
Assuming the determinant construction and Lemmas 4.2, 5.1, and 6.2 of
Banks, we obtain \(q<10^6\) for imaginary quadratic fields of class
number one. This improves the preceding cutoff \(5\cdot10^7\) by a
factor of 50. The key ingredient is a direct upper bound on the weighted
sum of split primes, obtained by counting primitive lattice points.
It replaces both the theta-integral zero construction and the
zero-dependent prime-bias estimate. Banks's determinant argument remains
an external input. This expanded version supplies the background,
determinant notation, and finite verification within one paper; no
earlier version is needed. The bound and analytic constants are unchanged.
No Heegner method or small-conductor zero computation is used.
\end{abstract}

\paragraph{Status.}
This is an AI-assisted research draft. Exact arithmetic checks and finite
computations are supplied, but an independent mathematical cross-check
of the v3 argument is pending. The completed check of v2 does not cover
this argument. The original v0.3.0 paper remains available unchanged.

\section{The problem, the result, and the route through the proof}

Tao's tenth-discriminant challenge \cite{tao} asks how large a further imaginary
quadratic field of class number one could be, using arguments related to
Siegel zeroes rather than Heegner's solution of the class-number-one
problem. Our aim is first to obtain an explicit analytic cutoff, and
then to check every smaller candidate by elementary integer arithmetic.
The two stages are distinct: the analytic argument supplies the bound
\(10^6\); the finite check leaves exactly the nine familiar fields.

Write \(K=\Q(\sqrt d)\), where \(d<0\) is squarefree. Its field
discriminant is the \emph{fundamental discriminant}
\[
 D=\begin{cases}
 d,&d\equiv1\pmod4,\\
 4d,&d\equiv2,3\pmod4.
 \end{cases}
\]
Put
\[
 q=|D|,\qquad K=\Q(\sqrt D),\qquad R=\sqrt q,\qquad L=\log q.
\]
Thus \(q\) is the absolute field discriminant; for example,
\(\Q(i)\) has \(q=4\). All logarithms are natural.
The class number \(h(K)\) counts ideal classes. We use the equivalence
\[
 h(K)=1
 \quad\Longleftrightarrow\quad
 \text{every nonzero integral ideal of \(\OK\) is principal}.
\]
Write \(\chi_D(n)=(D/n)\) for the quadratic Kronecker character of
conductor \(q\). The meaning of its prime values is explained below.

\begin{samepage}
\begin{theorem}\label{main}
Assume Banks's determinant construction and Lemmas 4.2, 5.1, and 6.2
in \cite{banks}, including the interpolation results used in that
construction. Then
\[
 h(K)=1\quad\Longrightarrow\quad q<1\,000\,000.
\]
Together with the finite verification in Section~\ref{finite}, this gives
\[
 h(K)=1\quad\Longleftrightarrow\quad
 q\in\{3,4,7,8,11,19,43,67,163\}.
\]
\end{theorem}
\end{samepage}

Here is the route through the proof.
\begin{enumerate}
\item Class number one turns prime ideals of norm \(p\) into algebraic
integers of norm \(p\). A split prime supplies two such lattice points
in the upper half-plane.
\item A geometric count, followed by a small inclusion--exclusion sieve,
bounds the weighted sum of split primes. Subtracting split and ramified
primes from an explicit total prime sum gives many inert primes.
\item Banks's construction supplies a nonzero integer \(J\). Its
absolute size has an upper bound, while each inert prime forces a large
prime-power divisor of \(J\). For \(q\ge10^6\), the lower bound exceeds
the upper bound.
\item Below \(10^6\), a separate exhaustive residue check excludes all
but nine discriminants. Reduced binary quadratic forms verify that
each survivor actually has class number one.
\end{enumerate}

The standard background is quadratic ideal theory and reduction of
positive binary quadratic forms; the facts needed from both are stated
in this paper. We also use the explicit prime estimates of
Rosser--Schoenfeld, stated in Section~\ref{prime-mass}.
Section~\ref{det-input} defines the determinant and states the exact
external bounds used from Banks. It does not reprove his interpolation
theorem. The present proof constructs no zero of an \(L\)-function and
does not use Banks's Theorem 1.1, Lemma 2.1, or Corollary 6.3.

\section{Quadratic integers, prime ideals, and the norm lattice}\label{background}

Embed \(K\) in \(\mathbb C\), choosing \(\sqrt D=i\sqrt q\).
The ring of integers has the explicit basis
\[
 \OK=\Z+\Z\omega,\qquad
 \omega=\begin{cases}
 (1+\sqrt D)/2,&D\text{ odd},\\
 \sqrt D/2,&D\text{ even}.
 \end{cases}
\]
In particular, \(\operatorname{Re}\omega=u\in\{0,1/2\}\) and
\(\operatorname{Im}\omega=R/2\). The field norm is
\(\operatorname{N}_{K/\Q}(\alpha)=\alpha\bar\alpha=|\alpha|^2\).
For integer coordinates \(a,b\), write
\[
 Q(a,b):=\operatorname{N}_{K/\Q}(a+b\omega)
        =(a+ub)^2+\frac q4b^2.
\]
It is a positive integer for \((a,b)\ne(0,0)\). An element is
nonrational exactly when \(b\ne0\), so every nonrational integer
has norm at least \(q/4\). An ordinary integer \(a\) has norm \(a^2\),
which cannot be prime.

For an integral ideal \(\mathfrak a\), its ideal norm is the index
\(\operatorname{N}\mathfrak a=[\OK:\mathfrak a]\).
For a principal ideal, this equals the field norm of a generator.
The three possibilities for a rational prime \(p\) are
\[
\begin{array}{c|c|c}
 \chi_D(p)&\text{factorization of }p\OK&\text{norms of prime ideals}\\ \hline
 1&\mathfrak p\,\bar{\mathfrak p},\ \mathfrak p\ne\bar{\mathfrak p}
      &p,\ p\\
 -1&p\OK\text{ is prime}&p^2\\
 0&\mathfrak p^2&p
\end{array}
\]
These are called split, inert, and ramified, respectively. Ramification
means \(p\mid D\). For odd \(p\), \(\chi_D(p)\) is zero if \(p\mid D\),
and otherwise is \(1\) or \(-1\) according as \(D\) is a square or a
nonsquare modulo \(p\). For \(p=2\), it is zero when \(D\) is even,
\(1\) when \(D\equiv1\pmod8\), and \(-1\) when \(D\equiv5\pmod8\).

\begin{lemma}[The small-prime obstruction]\label{small-prime}
If \(h(K)=1\), every prime \(p<q/4\) is inert.
\end{lemma}
\begin{proof}
A split or ramified prime has an ideal of norm \(p\). Under class
number one, that ideal is \((\alpha)\) with
\(\operatorname{N}_{K/\Q}(\alpha)=p\). The generator is nonrational,
so the norm formula gives \(p\ge q/4\), a contradiction.
\end{proof}

For \(q>4\), the only units are \(\pm1\). Thus an ideal with one
generator \(\alpha\) has exactly the two generators \(\alpha,-\alpha\).
If its norm is prime, neither is real, and exactly one has positive
imaginary part. Distinct split prime ideals therefore give exactly two
points with \(b>0\) and norm \(p\). This explains the factor \(1/2\)
in the split-prime estimate below.

\section{Counting lattice points and imposing primitivity}\label{counting}

For \(x\ge0\), define
\[
 A(x)=\#\{(a,b)\in\Z\times\Z_{\ge1}:Q(a,b)\le x\}.
\]
The region \(Q(a,b)\le x\) is an ellipse in the coordinate plane.
The restriction \(b>0\) counts one half of the nonrational points.

\begin{lemma}\label{lattice}
For every \(x\ge0\),
\[
 \frac{\pi x}{R}-2\sqrt x-\frac{2\sqrt x}{R}
 \le A(x)\le
 \frac{\pi x}{R}+\frac{2\sqrt x}{R}.
\]
\end{lemma}
\begin{proof}
Put \(B=2\sqrt x/R\) and
\(f(t)=2\sqrt{x-qt^2/4}\) on \([0,B]\).
For each integer \(1\le b\le B\), the allowed real values of \(a\)
form a closed interval of length \(f(b)\). Its integer count lies
between \(f(b)-1\) and \(f(b)+1\). Since \(f\) is decreasing,
\[
 \int_0^B f(t)\,dt-f(0)
 \le\sum_{b=1}^{\lfloor B\rfloor}f(b)
 \le\int_0^B f(t)\,dt.
\]
The integral is \(\pi x/R\), and there are at most \(B\) intervals.
This gives both bounds, including when \(B<1\).
\end{proof}

A pair \((a,b)\) is \emph{primitive} if \(\gcd(a,b)=1\). A common
divisor \(d>1\) would make \(d^2\) divide \(Q(a,b)\), so every pair of
prime norm is primitive. We only remove common factors among four
small primes. This leaves a superset of the primitive pairs, which
is sufficient for an upper bound. Set
\[
 P=2\cdot3\cdot5\cdot7=210,\qquad
 A_P(x)=\#\{(a,b):b\ge1,\ Q(a,b)\le x,\ \gcd(a,b,P)=1\}.
\]

\begin{lemma}\label{primitive}
For every \(x\ge0\),
\[
 A_P(x)\le\frac{\pi\rho}{R}x+B_R\sqrt x,
 \qquad
 \rho=\frac{768}{1225},\quad
 B_R=\frac{88}{35}+\frac{192}{35R}.
\]
\end{lemma}
\begin{proof}
Write \(\mu(d)=(-1)^k\) when \(d\mid P\) has \(k\) prime factors.
The pairs whose two coordinates are divisible by \(d\) correspond
to the pairs counted by \(A(x/d^2)\). Inclusion--exclusion therefore gives
\[
 A_P(x)=\sum_{d\mid P}\mu(d)A(x/d^2).
\]
For positive coefficients use the upper bound of Lemma~\ref{lattice};
for negative coefficients use its lower bound. The area coefficient is
\[
 \sum_{d\mid P}\frac{\mu(d)}{d^2}
   =\prod_{p\mid P}(1-p^{-2})=\rho.
\]
Write
\[
 T=\sum_{d\mid P}\frac{|\mu(d)|}{d}=\frac{96}{35},\qquad
 U=\sum_{d\mid P}\frac{\mu(d)}d=\frac8{35}.
\]
The boundary coefficient is
\(2\sum_{\mu(d)=-1}d^{-1}+2T/R=T-U+2T/R=B_R\).
\end{proof}

\section{Direct bounds on split and ramified primes}

Assume \(h(K)=1\) and \(q\ge10^6\), and put \(a_0=q/4\).
If \(p\) splits, its two prime ideals of norm \(p\) are principal.
As explained in Section~\ref{background}, there are exactly four generators
of norm \(p\). They give two distinct pairs with \(b>0\).
They are nonrational, and each satisfies \(\gcd(a,b,P)=1\), since
\(d\mid a,b\) would force \(d^2\mid Q(a,b)=p\).
Consequently,
\begin{equation}\label{split-count}
 S_+(X):=\sum_{\substack{p\le X\\\chi_D(p)=1}}\frac{\log p}{p}
 \le\frac12\sum_{\substack{b\ge1,\ Q(a,b)\le X\\\gcd(a,b,P)=1}}
                        \frac{\log Q(a,b)}{Q(a,b)}.
\end{equation}

All norms in this sum are at least \(a_0\).
For \(f(t)=(\log t)/t\), we have
\(-f'(t)=(\log t-1)/t^2>0\) on \([a_0,\infty)\).
Partial summation expresses the weighted count as its endpoint term
plus \(-\int A_P(t)f'(t)\,dt\); the lower endpoint uses
\(A_P(a_0^-)=0\), so points of norm exactly \(a_0\) are included.
For \(X\ge a_0\), inserting Lemma~\ref{primitive} gives
\begin{align}
 S_+(X)
 &\le\frac12\left\{
 A_P(X)\frac{\log X}{X}
 +\int_{a_0}^X A_P(t)\frac{\log t-1}{t^2}\,dt\right\} \notag\\
 &\le\frac{\pi\rho}{4R}
       \bigl((\log X)^2-(\log a_0)^2+2\log a_0\bigr)
   +\frac{B_R(\log a_0+1)}{\sqrt{a_0}}.\label{split-precise}
\end{align}
In the last line we dropped the nonpositive boundary term
\(-B_R(\log X+2)/(2\sqrt X)\).
Here \(\log a_0>2\), and \(\log a_0+1<L\) because \(\log4>1\).
Also, for \(R\ge1000\),
\[
 \frac{\pi\rho}{4}<\frac{493}{1000},\qquad B_R<\frac{63}{25}.
\]
Thus the required direct split-prime bound is
\begin{equation}\label{split-final}
 S_+(X)<\frac{493}{1000R}(\log X)^2+\frac{126L}{25R}
             \qquad(X\ge q/4).
\end{equation}

A ramified prime likewise has a principal prime ideal of norm \(p\).
Its generator is nonrational, so \(p\ge q/4\). Hence
\begin{equation}\label{ramified}
 S_0(X):=\sum_{\substack{p\le X\\p\mid q}}\frac{\log p}{p}
 \le\frac4q\sum_{p\mid q}\log p\le\frac{4L}{q}.
\end{equation}
Neither this estimate nor the split-prime count requires a separate
reduction to prime discriminants.

\section{Obtaining the inert-prime mass}\label{prime-mass}

For any \(X>H\ge319\), the explicit prime estimates we use are
\begin{equation}\label{mertens}
 \sum_{H<p\le X}\frac{\log p}{p}
 >\log X-\log H-\frac1{\log H},
 \qquad
 \vartheta(X):=\sum_{p\le X}\log p<2X.
\end{equation}
For completeness, Rosser--Schoenfeld's Theorem 6 \cite{rs} gives a constant \(C\)
such that
\[
 \left|\sum_{p\le x}\frac{\log p}{p}-\log x-C\right|
 <\frac1{2\log x}\qquad(x\ge319).
\]
Subtract the upper bound at \(H\) from the lower bound at \(X\).
Since \((2\log X)^{-1}<(2\log H)^{-1}\), this proves the first
inequality in \eqref{mertens}. Their Theorem 9 gives the stronger
\(\vartheta(x)<1.01624x\) for \(x>0\), so also the second.

Every prime is split, inert, or ramified. Removing the full split and
ramified sums up to \(X\) removes at least their contributions in
\((H,X]\). Thus
\[
 \sum_{\substack{H<p\le X\\\chi_D(p)=-1}}\frac{\log p}{p}
 \ge \sum_{H<p\le X}\frac{\log p}{p}-S_+(X)-S_0(X).
\]
Combining \eqref{mertens}, \eqref{split-final}, and \eqref{ramified},
for \(X>H\ge319\) and \(X\ge q/4\) we obtain
\begin{equation}\label{inert}
 \sum_{\substack{H<p\le X\\\chi_D(p)=-1}}\frac{\log p}{p}
 >\log X-\log H-\frac1{\log H}
  -\frac{493}{1000R}(\log X)^2
  -\left(\frac{126}{25R}+\frac4q\right)L.
\end{equation}

\section{The determinant and the external results used}\label{det-input}

We now define the objects whose estimates are taken from Banks. This
section is an explicit statement of that input, not a new proof of his
interpolation result. Our notation \(K\) denotes the quadratic field;
the biquadratic field called \(K\) in Banks is denoted by \(E\) here.

\subsection{The points, the rows, and the nonzero integer}

Let \(d<0\) be the squarefree integer from Section 1, and set
\[
 \eta=\sqrt d,\qquad \xi=\sqrt2,\qquad
 E=\Q(\eta,\xi),\qquad \mathcal O=\Z[\eta,\xi].
\]
The field \(E\) has degree four over \(\Q\). Its four automorphisms
are \(1,\sigma,\tau,\sigma\tau\), where \(\sigma\) changes the sign
of \(\eta\) and \(\tau\) changes the sign of \(\xi\). Elements of
\(\mathcal O\) are algebraic integers, so their field norms are integers.

Choose integers \(1\le H\le N\), and put \(M=N^{4/3}\). The real number
\(M\) need not be an integer. The column set is
\(\mathcal I=\{0,\ldots,N-1\}^4\), of cardinality \(N^4\). For each
\(\boldsymbol n=(n_1,n_2,n_3,n_4)\in\mathcal I\), set
\[
 \theta_{\boldsymbol n}
 =n_1+n_2\eta+n_3\xi+n_4\eta\xi.
\]
For \(\boldsymbol\alpha=(\alpha_1,\alpha_2,\alpha_3)
\in\Z_{\ge0}^3\), form the row
\[
 \boldsymbol v_{\boldsymbol\alpha}
 =\bigl(\theta_{\boldsymbol n}^{\alpha_1}
       \sigma(\theta_{\boldsymbol n})^{\alpha_2}
       \sigma\tau(\theta_{\boldsymbol n})^{\alpha_3}
   \bigr)_{\boldsymbol n\in\mathcal I},
 \qquad
 w(\boldsymbol\alpha)=\alpha_1+H\alpha_2+H\alpha_3.
\]
Here \(z^0=1\), including at \(z=0\). Order the indices by increasing
weight, breaking ties by any fixed rule, and retain a row if it is
linearly independent over \(E\) of the earlier retained rows.
Let \(\mathcal A\) be the retained index set.

The interpolation and construction results in Banks's Sections 3--4
give \(|\mathcal A|=N^4\) and a basis of the \(N^4\)-dimensional row
space. Consequently, the following determinant and integer are nonzero:
\[
 \Delta=\det(\boldsymbol v_{\boldsymbol\alpha})_{\boldsymbol\alpha\in\mathcal A}
       \in\mathcal O\setminus\{0\},\qquad
 J=\operatorname{N}_{E/\Q}(\Delta)
   =\Delta\,\sigma(\Delta)\,\tau(\Delta)\,\sigma\tau(\Delta)
       \in\Z\setminus\{0\}.
\]
Define the degree sums
\[
 S_1=\sum_{\boldsymbol\alpha\in\mathcal A}\alpha_1,\qquad
 S_{23}=\sum_{\boldsymbol\alpha\in\mathcal A}(\alpha_2+\alpha_3).
\]
These are the quantities denoted by \(S_1,S_{23}\) below.
The argument uses their bounds; it never numerically constructs this
enormous determinant.

\subsection{The three estimates}

For \(q\ge9\) and \(H^{2/3}N^{4/3}\ge2\cdot97^2\), the precise
consequences of the cited lemmas are as follows.
\begin{enumerate}
\item \emph{Degree estimates (Lemma 4.2):}
\begin{equation}\label{banks-degree}
 S_1\ge\frac{N^4H^{2/3}M}{4\cdot97^2}>0,
 \qquad
 S_{23}\le192N^4H^{-1/3}M.
\end{equation}
\item \emph{Absolute-size estimate (Lemma 5.1):}
\begin{equation}\label{banks-arch}
 \frac14\log|J|
 \le (S_1+S_{23})\left(\log N+\frac L2+\log8\right)
                       +2N^4\log N.
\end{equation}
This bounds the determinant at its four complex embeddings.
\item \emph{Prime-power divisibility (Lemma 6.2):}
for every odd prime \(p>H\) with \(\chi_D(p)=-1\),
\begin{equation}\label{banks-div}
 p^{4E_p}\mid J,\qquad
 E_p=\sum_{\boldsymbol\alpha\in\mathcal A}
           \left\lfloor\frac{\alpha_1}{p}\right\rfloor.
\end{equation}
The factor four comes from taking the norm from \(E\): a factor
\(p^{E_p}\) of \(\Delta\) contributes \(p^{4E_p}\) to its norm.
\end{enumerate}

The mechanism in the last estimate is the Frobenius congruence
\(\theta^p\equiv g_p(\theta)\pmod{p\mathcal O}\), where \(g_p\)
is \(\sigma\) or \(\sigma\tau\). In a row, exchanging \(p\) powers
of the first variable for one power of either other variable decreases
the weight when \(p>H\). Banks uses the resulting row operations to
extract the stated prime powers. No hypothesis about a real zero is
present in \eqref{banks-degree}--\eqref{banks-div}.

\section{Applying the determinant estimates}\label{comparison}

Take
\[
 H=10^9,\qquad N=q^{13},\qquad M=N^{4/3},\qquad
 z=\frac{52}{3},\qquad \log M=zL.
\]
For \(q\ge10^6\), Banks's hypotheses
\(q\ge9\), \(1\le H\le N\), and
\(H^{2/3}N^{4/3}\ge2\cdot97^2\) hold, as do \(M>H\) and \(M\ge q/4\).
All definitions and estimates of Section~\ref{det-input} therefore apply.
Dividing the degree bounds \eqref{banks-degree} gives
\[
 \frac{S_{23}}{S_1}\le\frac{768\cdot97^2}{10^9}
                  <r:=\frac{73}{10000},\qquad
 \frac{N^4M}{S_1}\le\frac{4\cdot97^2}{10^6}
                  <c:=\frac{19}{500}.
\]
Divide \eqref{banks-arch} by \(S_1\).
Since \(N^4/S_1<c/M\), it gives
\begin{equation}\label{arch}
 \frac{\log|J|}{4S_1}
 \le(1+r)\left(\log N+\frac L2+\log8\right)
                         +\frac{2c\log N}{M}.
\end{equation}

There are \(N^4\) retained indices. The elementary bound
\(\lfloor t\rfloor\ge t-1\) therefore gives
\[
 E_p\ge\frac{S_1}{p}-N^4.
\]
The powers in \eqref{banks-div} for different primes are coprime, so
their product divides the nonzero integer \(J\). Hence
\begin{align*}
 \frac{\log|J|}{4S_1}
 &\ge\sum_{\substack{H<p\le M\\\chi_D(p)=-1}}
                       \frac{E_p}{S_1}\log p\\
 &\ge\sum_{\substack{H<p\le M\\\chi_D(p)=-1}}\frac{\log p}{p}
       -\frac{N^4}{S_1}
        \sum_{\substack{H<p\le M\\\chi_D(p)=-1}}\log p.
\end{align*}
The last sum is at most \(\vartheta(M)<2M\), and
\(N^4M/S_1<c\). Consequently,
\begin{equation}\label{padic}
 \frac{\log|J|}{4S_1}
 \ge\sum_{\substack{H<p\le M\\\chi_D(p)=-1}}\frac{\log p}{p}-2c.
\end{equation}
This derives the lower bound from prime-power divisibility itself.

Compare \eqref{arch} with \eqref{padic} and insert \eqref{inert}
at \(X=M\). The leading lower term is \(\log M\), whereas the leading
upper term is \((1+r)\log N=\tfrac34(1+r)\log M\).
Moving the other terms to the right, using \(\log N=13L\), and
dividing by \(L\) gives the necessary inequality
\begin{align}
 \left(\frac14-\frac{3r}{4}\right)z
 &\le \frac{493}{1000}z^2\frac{L}{R}
       +\frac{1+r}{2}+\frac{126}{25R}+\frac4q \notag\\
 &\quad+\frac{\log H+(1+r)\log8+1/\log H+2c}{L}
       +\frac{26c}{M}.\label{necessary}
\end{align}

\section{A uniform rational contradiction}\label{contradiction}

The derivative of \((\log q)/\sqrt q\) is
\(q^{-3/2}(1-\tfrac12\log q)\), so the function decreases for \(q>e^2\).
Using \(q\ge10^6\), we have the uniform estimates
\[
 \begin{gathered}
 \frac LR<\frac{13.816}{1000},\qquad L>13.815,\qquad R\ge1000,\qquad
 \frac{\log H}{L}\le\frac32,\\
 \log8<2.08,\qquad \frac1{\log H}<0.049,\qquad
 \frac{26c}{M}<10^{-6}.
 \end{gathered}
\]
The last inequality follows already from \(M\ge10^6\) and
\(26c=0.988<1\). The logarithmic estimates are certified by rational
Taylor bounds for the exponential function in the accompanying program.

The left side of \eqref{necessary} equals
\[
 \left(\frac14-\frac34\frac{73}{10000}\right)\frac{52}{3}
                  =\frac{127153}{30000}>4.238.
\]
Its right side is strictly less than the rational number
\begin{align*}
 \mathcal R={}&\frac{493}{1000}\left(\frac{52}{3}\right)^2
                       \frac{13.816}{1000}
             +\frac{1.0073}{2}+\frac32\\
 &+\frac{1.0073(2.08)+0.049+0.076}{13.815}
             +\frac{5.04}{1000}+\frac4{10^6}+\frac1{10^6}
             <4.216.
\end{align*}
Every terminating decimal here denotes its exact rational value.
More precisely, the certificate checks
\[
 \frac{127153}{30000}-\mathcal R
   =\frac{867907363}{38375000000}>0.022.
\]
Every term has been bounded uniformly for all \(q\ge10^6\), rather
than only at the endpoint. This contradicts \eqref{necessary} and
proves the analytic cutoff in Theorem~\ref{main}.

\section{The complete finite verification below the cutoff}\label{finite}

We now justify the finite classification within this paper. It uses
Lemma~\ref{small-prime}, not the determinant estimates. The finite
discriminant search here is also distinct from the sieve by \(2,3,5,7\)
in Section~\ref{counting}, which counts lattice points for one field.

\subsection{Why these are all the candidates}

If \(q>8\) and \(h(K)=1\), Lemma~\ref{small-prime} applies to \(p=2\).
The splitting rule in Section~\ref{background} then gives
\[
 D\equiv5\pmod8,\qquad q\equiv3\pmod8.
\]
The fundamental negative discriminants with \(q\le8\) are exactly
\(-3,-4,-7,-8\), by the definition in Section 1. Thus, after the
analytic cutoff, every additional candidate lies in the finite set
\[
 11\le q<1\,000\,000,\qquad q\equiv3\pmod8.
\]
For each odd prime \(p\) with \(4p<q\), the integer \(-q\) must be a
nonzero nonsquare modulo \(p\). We exclude \(q\) if it fails even one
of these necessary conditions.

It is sufficient to use only primes \(p\le109\). The direct scan checks
all \(124\,999\) integers in the displayed progression. It deliberately
includes nonfundamental integers: omitting the squarefreeness test
enlarges the candidate set and cannot discard a genuine field.
Let \(T_p=\{a^2\bmod p:0\le a<p\}\), including zero. The complete
exclusion rule can be written as the following Python pseudocode:
\begin{verbatim}
primes = [odd primes at most 109, generated by trial division]
squares = {p: {a*a % p for a in range(p)} for p in primes}
survivors = []
for q in range(11, 1_000_000, 8):
    if all((-q) % p not in squares[p]
           for p in primes if 4*p < q):
        survivors.append(q)
\end{verbatim}
The supplied executable program \texttt{verify\_finite\_range.py}
implements this scan using only integer arithmetic. It gives exactly
\[
 \texttt{survivors}=[11,19,43,67,163].
\]
Its output records how many candidates each first rejecting prime
eliminates; those counts sum to \(124\,994\). Adding the five survivors
accounts for every scanned integer.

\subsection{Why each surviving field has class number one}

Passing the residue tests is only a necessary condition. To verify the
survivors, use the classical correspondence between ideal classes of
\(\OK\) and proper equivalence classes of primitive positive definite
integral binary quadratic forms
\[
 f(x,y)=ax^2+bxy+cy^2,\qquad b^2-4ac=D=-q.
\]
Here primitive means \(\gcd(a,b,c)=1\), and proper equivalence means
an integral change of variables of determinant one. For example, the
form determines the ideal class represented by
\(a\Z+\frac{-b+\sqrt D}{2}\Z\).
The reduction theorem says that each class has exactly one representative
with
\[
 |b|\le a\le c,\qquad
 b\ge0\ \text{if }|b|=a\text{ or }a=c.
\]
Thus the number of these reduced forms is \(h(K)\).
The inequalities imply
\[
 q=4ac-b^2\ge3a^2,\qquad
 1\le a\le\left\lfloor\sqrt{q/3}\right\rfloor.
\]
This makes the check exhaustive: enumerate these \(a\), then
\(-a\le b\le a\), and retain
\(c=(b^2+q)/(4a)\) exactly when it is an integer with \(c\ge a\),
the gcd condition holds, and the boundary convention holds.

The program checks that every survivor is a fundamental discriminant
and enumerates its reduced forms. The complete result is
\[
\begin{array}{r|c@{\qquad}r|c@{\qquad}r|c}
 q&(a,b,c)&q&(a,b,c)&q&(a,b,c)\\ \hline
 3&(1,1,1)&8&(1,0,2)&43&(1,1,11)\\
 4&(1,0,1)&11&(1,1,3)&67&(1,1,17)\\
 7&(1,1,2)&19&(1,1,5)&163&(1,1,41)
\end{array}
\]
Each field has exactly one reduced form. This proves the remaining
assertion of Theorem~\ref{main}, subject to the explicitly stated
analytic inputs and the reproducible finite computation.

\subsection{Relation to the larger existing certificate}

The bundle also retains \texttt{verify\_sieve.cpp} and
\texttt{verify\_independent.py}, with their original outputs. They check
the larger interval \(q<10^{10}\), which contains the present one.
For \(q\ge437\), they first combine the conditions for
\(2,3,5,7,11,13,17\) into \(1440\) residue classes modulo
\(8\cdot3\cdot5\cdot7\cdot11\cdot13\cdot17=2042040\).
They then test primes from \(19\) through \(109\). All are eligible
there because \(4\cdot109<437\). Of \(7\,051\,769\) candidates in that
range, none survive. For \(11\le q<437\), they apply only the eligible
conditions \(4p<q\), leaving the same five values.
The C++ program builds residue classes by the Chinese remainder theorem;
the Python program scans the initial classes and uses periodic bitsets.
This is supporting evidence. The direct scan just described suffices
to reproduce the finite part of the present proof.

\section{Reproduction, verification scope, and version history}\label{reproduce}

Run these commands from this version's directory:
\begin{verbatim}
python3 verify_direct_prime.py --sanity
python3 verify_finite_range.py
\end{verbatim}
Both use only the Python standard library.
The first records its output in \texttt{checks.txt} and passes 26 exact
arithmetic checks and 22 finite sanity cases.
The second checks the entire remaining candidate progression and the
nine reduced-form counts; its output is in
\texttt{finite\_range\_checks.txt}.

\subsection{What the arithmetic certificate establishes}

The 26 checks verify the sieve constants, determinant parameters,
logarithmic estimates, and final rational contradiction. For
\(x\ge0\) and an integer \(n\) with \(n+2>x\), put
\(T_n(x)=\sum_{k=0}^n x^k/k!\). The bound used in the program is
\[
 T_n(x)\le e^x\le T_n(x)+
 \frac{x^{n+1}}{(n+1)!}\frac1{1-x/(n+2)}.
\]
Indeed, after the first omitted term, each ratio of successive terms
is at most \(x/(n+2)<1\), so a geometric series bounds the tail.
All quantities are exact fractions. With \(n=160\), the program
certifies, for example, \(e^{13.815}<10^6<e^{13.816}\), which proves
the two logarithmic bounds. The classical inequality \(\pi<22/7\)
is used to check the coefficient \(0.493\).

The 22 sanity cases comprise 15 finite lattice comparisons and seven
prime-norm multiplicity checks. The lattice examples compare direct
enumeration with row counting, inclusion--exclusion, and the geometric
bounds. The prime checks cover all primes at most \(2000\) for each
of \(q=7,8,11,19,43,67,163\). These examples test the formulas and
implementation; they do not prove the general counting lemmas.

\subsection{What changed, and what remains external}

The original v3 cutoff \(q<10^6\) \cite{v3}, all its analytic constants, and its
26-check arithmetic certificate are unchanged in v0.3.1. This revision
adds the setup and roadmap, defines the determinant objects and states
their input bounds, expands the intermediate deductions, and includes
the finite classification and a direct verifier for the exact remaining
range. The preceding v2 cutoff was \(q<5\cdot10^7\); its theta-integral
route is historical context and is not needed here.

Banks's interpolation and determinant results, the Rosser--Schoenfeld
prime estimates, and the stated classical quadratic theory remain
mathematical inputs. The programs do not audit those results or
formalize the general analytic deductions. No independent mathematical
cross-check of v3 or Lean formalization is recorded here.
Codex generated the v3 derivation, this expanded exposition, and the
verification programs. The completed independent check of v2 remains
scoped to v2. Source URLs and snapshot hashes are in
\texttt{sources.json}; detailed review status is in \texttt{REVIEW.md}.

\begin{thebibliography}{9}
\bibitem{banks}
W.~D.~Banks, \emph{On the OpenAI proof that Landau--Siegel zeros
do not exist}, \href{https://arxiv.org/abs/2610.15483v1}
{arXiv:2610.15483v1}, especially the determinant construction in
Sections 3--4 and Lemmas 4.2, 5.1, and 6.2.
\bibitem{rs}
J.~B.~Rosser and L.~Schoenfeld,
\emph{Approximate formulas for some functions of prime numbers},
Illinois J. Math. \textbf{6} (1962), 64--94, Theorems 6 and 9.
\href{https://doi.org/10.1215/ijm/1255631807}{doi:10.1215/ijm/1255631807}.
\bibitem{tao}
T.~Tao, \emph{An experimental AI-math challenge: using tenth-discriminant
bounds to explore arguments ruling out Siegel zeroes},
\href{https://terrytao.wordpress.com/2026/10/20/an-experimental-ai-math-challenge-using-tenth-discriminant-bounds-to-explore-arguments-ruling-out-siegel-zeroes/}
{challenge and discussion}, 20 October 2026.
\bibitem{v3}
\emph{A direct split-prime bound for the tenth discriminant},
\href{https://github.com/twinheterodox/tenth-discriminant/releases/tag/v0.3.0}
{original v3 release v0.3.0}, 20 October 2026.
\end{thebibliography}
\end{document}
